\documentclass[11pt]{article}
\usepackage{hyperref}
\usepackage{cmap}
\usepackage[utf8]{inputenc}
\usepackage[T1]{fontenc}
\usepackage[left=2.25cm,top=1.8cm,right=2.25cm,bottom=1.8cm]{geometry}
\IfFileExists{glyphtounicode.tex}{\input{glyphtounicode}\pdfgentounicode=1}{}
\hypersetup{
    hidelinks,
    pdfauthor={Kevin Bravo},
    pdftitle={Kevin Bravo - Carta de presentacion - Proyecto Spider-Man Brand New Day},
    pdfsubject={Carta de presentacion para el proyecto frontend de InfoJobs y Midudev},
    pdfkeywords={Kevin Bravo, carta de presentacion, frontend developer, InfoJobs, Midudev, Spider-Man Brand New Day, React, Astro, SvelteKit, TypeScript}
}
\setlength{\parindent}{0pt}
\setlength{\parskip}{9pt}
\pagestyle{empty}

\begin{document}

{\LARGE \textbf{Kevin Bravo}}\\[4pt]
Frontend Developer | React, Astro, SvelteKit, TypeScript\\[8pt]
\href{mailto:0bravokevin@gmail.com}{0bravokevin@gmail.com}
\quad\textbullet\quad
\href{https://kevinbravo.com}{kevinbravo.com}
\quad\textbullet\quad
\href{https://github.com/0bkevin}{github.com/0bkevin}

\vspace{8pt}
\hrule
\vspace{14pt}

\textbf{Equipo de InfoJobs y Midudev:}

Quiero participar en el desarrollo de la landing de \textit{SPIDER-MAN: Brand New Day}. El proyecto combina tres aspectos en los que tengo experiencia directa: crear interfaces cuidadas, optimizar productos web para una navegación rápida y entregar funcionalidades completas en plazos cortos.

Actualmente desarrollo RegistroCiberVE, un registro público bilingüe construido con Svelte 5, SvelteKit, TypeScript y Tailwind CSS. Diseñé su landing, directorios, búsqueda, filtros, perfiles y visualizaciones responsivas. También implementé SSR, SEO técnico, sitemap dinámico, contenido para compartir y el despliegue en Cloudflare Workers. Este trabajo me ha exigido equilibrar identidad visual, claridad, rendimiento y comportamiento en distintos tamaños de pantalla.

En Free2Z lideré la migración de un producto activo de React a SvelteKit. Reorganicé la arquitectura frontend y mejoré SSR, mantenibilidad, estado y rendimiento sin interrumpir el servicio. Además, integré APIs en Python y Django, video en vivo con Dyte y herramientas de publicación. Como freelance he desarrollado sitios con Astro, TypeScript y Tailwind CSS, trabajando desde la definición de la interfaz hasta la publicación y los ajustes con el cliente.

Puedo aportar autonomía para resolver problemas, criterio visual y experiencia llevando interfaces desde una idea hasta producción. También disfruto el trabajo colaborativo. Como organizador del hub de Caracas del Dev3pack Global Hackathon coordiné participantes, mentores y aliados, y ayudé a equipos a definir el alcance y entregar sus proyectos dentro del plazo.

Me atrae el formato intensivo de dos semanas y la oportunidad de construir una experiencia que recibirá mucha atención desde su lanzamiento. Estaría encantado de conversar sobre el enfoque técnico y creativo del proyecto.

Saludos,\\[3pt]
\textbf{Kevin Bravo}

\end{document}
